\chapter*{Introduction}
\addcontentsline{toc}{chapter}{Introduction}

The simplest stochastic differential equation involving the local time of the
unknown process is
\[
	X_t = x + W_t + \beta\hat{L}^0_t(X),
\]
where $W$ is a Brownian motion and $\hat{L}^0(X)$ is the symmetric local time
of $X$ at $0$. Its solution, the skew Brownian motion, behaves like a Brownian
motion away from $0$, but each excursion away from $0$ is positive with
probability $\frac{1+\beta}{2}$ and negative with probability
$\frac{1-\beta}{2}$. Its transition semigroup solves the heat equation on each
side of $0$, together with a \emph{transmission condition} at $0$:
\[
	(1+\beta)\frac{\partial u}{\partial x}(t,0^+) = (1-\beta)\frac{\partial u}{\partial x}(t,0^-).
\]
Equations of this kind model diffusion in media with an interface, where the
coefficients of the equation are discontinuous. This thesis studies
stochastic differential equations involving local times, and the parabolic
partial differential equations with transmission condition associated with
them.

\begin{description}
	\item[Chapter~1] recalls the construction and the main properties of the
	      local time of a continuous semimartingale: Tanaka's formula, the
	      Itô--Tanaka formula, the occupation times formula, the regularity of
	      local times in the space variable, the approximation by down-crossings,
	      and the symmetric local time, which is used throughout the rest of the
	      thesis.
	\item[Chapter~2] studies the time-inhomogeneous equation
	      \[
		      X_t = x + B_t + \int_0^t \alpha(s)\,dL_s ,
	      \]
	      where $L$ is the local time of $X$ at $0$. Following the method of Le
	      Gall, we prove pathwise uniqueness when $\alpha\leqslant\frac12$, and
	      extend the result to equations of the form
	      \[
		      X_t = x + \int_0^t \sigma(X_s)\,dB_s + \int_0^t b(X_s)\,ds + \int_0^t \alpha(s)\,dL_s .
	      \]
	\item[Chapter~3] follows \citet{LeGall1984}. It first gives pathwise
	      uniqueness results for equations $dX_t = \sigma(t,X_t)\,dB_t + b(t,X_t)\,dt$
	      with non-Lipschitz coefficients. These are then used to study the
	      singular equation
	      \[
		      dX_t = \phi(X_t)\,dB_t + \int_{\mathbb{R}} \nu(da)\,d_t\hat{L}^a_t(X),
	      \]
	      which involves the local times of the unknown process at all levels.
	      The chapter ends with limit theorems and comparison theorems.
	\item[Chapter~4] recalls, following \citet{Durrett1996}, how Brownian
	      motion gives the solutions of the heat equation and of the
	      Feynman--Kac equation.
	\item[Chapter~5] applies these ideas to the skew Brownian motion. We
	      construct it by random flipping of the excursions of a reflected
	      Brownian motion, compute the joint law of the process and its local
	      time, deduce its transition density, and derive the partial
	      differential equation with transmission condition that it solves.
\end{description}

\paragraph{About this version.} This thesis was written in 2018 for the
Master's degree M2PA at Cadi Ayyad University. The present version was
revised in 2026: the mathematical errors of the original have been
corrected, several proofs have been completed, and references have been
added. The list of corrections is kept with the source files.
